\chapter{The Tropics: ITCZ, Monsoons, and Variability}
\label{chap:ch19_tropics_monsoon}

The tropics are the engine room of the climate system. Solar heating peaks near the equator,
moist convection lofts that heat into the upper troposphere, and the resulting Hadley
circulation exports it poleward. The narrow, deep-convecting rain band where the trade winds of
the two hemispheres meet --- the \emph{Intertropical Convergence Zone} (ITCZ) --- is the single
most prominent feature of the observed precipitation field, and one of the hardest for a coarse
model to capture. This chapter evaluates how \model{} represents the tropical mean state: the
ITCZ and its precipitation, what role the multi-level moist atmosphere plays in producing it,
and --- honestly --- where the representation remains weak, including the near-absence of
interannual (ENSO-like) variability and the limited monsoon at this resolution.

This is a Part~IV evaluation chapter. We first define the phenomenon and the physics that should
produce it, recall how the relevant model components implement that physics
(Chapters~\ref{chap:ch09_atmosphere_multilevel} and \ref{chap:ch10_atmosphere_singlelevel}),
describe the diagnostic and the skill metric, then present the validation figures and a
qualitative assessment with its limitations.

\section{The phenomenon: deep convection and the ITCZ}

In the time mean, the trade-wind boundary layers of the Northern and Southern Hemispheres
converge in the deep tropics, forcing ascent, condensation, and a band of intense rainfall. The
ascending branch of the Hadley cell sits over this band; the descending branches dry the
subtropics, producing the desert latitudes near $\pm25^\circ$. The position and strength of the
ITCZ are set by the cross-equatorial energy transport and by the underlying sea-surface
temperature (SST), since deep convection switches on only above an SST threshold of roughly
$26$--$27^\circ$C \citep{vallis2017}. A faithful tropical climate therefore requires two
ingredients simultaneously: a moist atmosphere that can convect and rain, and an ocean whose
tropical SST is warm enough to trigger that convection.

\section{Why vertical structure matters: single level vs.\ multi-level}

The distinction between \model{}'s two atmospheres is sharpest in the tropics. The
\textbf{single-level (barotropic) atmosphere} of Chapter~\ref{chap:ch10_atmosphere_singlelevel}
has no vertical coordinate: it cannot represent the ascent--condensation--latent-heating column
that \emph{is} deep tropical convection. With no resolved Hadley overturning to concentrate
moisture at the equator, its precipitation is too dry and rains in the wrong place --- in the
subtropics and mid-latitudes rather than on the equator --- and its precipitation pattern
correlation against observations is near zero. This is a structural limitation, not a tuning
failure: a single level simply has nowhere to put the vertical motion that makes an ITCZ.

The \textbf{multi-level spectral atmosphere} built on the \dino{} dycore
(Chapter~\ref{chap:ch09_atmosphere_multilevel}) does have the necessary vertical structure. It
carries prognostic specific humidity as an advected tracer, evaporates moisture from the surface
with a bulk formula \citep{largeyeager2004}, and rains out supersaturation through a large-scale condensation closure
that latent-heats the column. As the resolved Hadley circulation spins up, that closure
concentrates rainfall in the deep tropics and dries the subtropics --- a real ITCZ forming
\emph{from the resolved circulation}, with no explicit convection scheme. \model{} therefore
follows the same lineage as NeuralGCM, whose differentiable spectral core \dino{} is
\citep{kochkov2024neuralgcm}. This is the central evaluation message of the chapter: the
multi-level atmosphere produces an equatorial rain band that the single level cannot.

\section{Diagnostic and metric}

The precipitation diagnostic is the column-integrated condensation rate, $P$ (in mm/day),
output by the multi-level atmosphere's moisture physics and regridded from the Gaussian
transform grid to the model's lat--lon grid. We assess it three ways against the GPCP-class
observational climatology bundled with the validation framework
(\code{chronos\_esm/validation/}):

\begin{description}[leftmargin=0pt,style=nextline]
  \item[Spatial map and bias map] $P$, the observed climatology, and their difference
    $P_{\mathrm{model}}-P_{\mathrm{obs}}$, to localize where the model is too wet or too dry.
  \item[Zonal mean] the latitude profile of $\langle P\rangle_\lambda$, model versus
    observations, which isolates the ITCZ peak, its width, and the subtropical minima.
  \item[Scalar scorecard] area-weighted bias, RMSE, pattern correlation, and standard-deviation
    ratio (a perfect model scores bias $0$, corr $1$, std ratio $1$).
\end{description}

\codref{experiments/make\_readme\_figures.py}{precipitation map, zonal mean, and scorecard}

\section{The current validation dashboard}

Figures~\ref{fig:dino_circ}--\ref{fig:precip_bias} are the present tropical validation
dashboard, generated from a short SST-forced spin-up of the multi-level atmosphere on the WOA18
SST climatology \citep{locarnini2018woa}. The quantitative scorecard is refreshed from the latest control run; a 30-year
time mean of the coupled \dino{}\,$\leftrightarrow$\,ocean control is being generated to replace
the short-spin-up numbers quoted below with an equilibrated climatology.

\begin{figure}[h]
\centering
\includegraphics[width=\linewidth]{dino_circulation}
\caption{Multi-level \dino{} atmosphere forced by the WOA SST (time mean of a $120$-day
spin-up). \emph{Left:} the zonal-mean zonal wind, with the twin upper-tropospheric
eddy-driven baroclinic jets ($\sim27$~m/s) that the single level cannot generate. \emph{Centre:}
surface zonal wind, showing trade easterlies straddling the equator. \emph{Right:} precipitation
(mm/day) --- the tropical ITCZ rain band, the feature of interest in this chapter. Regenerate
with \code{python experiments/dino\_circulation\_figure.py}.}
\label{fig:dino_circ}
\end{figure}

The right-hand panel of Figure~\ref{fig:dino_circ} is the headline result: an equatorial band of
enhanced precipitation, organized by the resolved Hadley ascent, with dry subtropical flanks.
The accompanying trade easterlies (centre panel) converging toward that band are the surface
signature of the same circulation. A single-level atmosphere produces neither.

\begin{figure}[h]
\centering
\includegraphics[width=0.72\linewidth]{zonal_precip}
\caption{Zonal-mean precipitation versus latitude: \model{}'s multi-level atmosphere (solid) and
the observed climatology (dashed). The model reproduces a distinct equatorial peak near
$5$--$15^\circ$N --- a genuine ITCZ --- but it is too dry: the peak undershoots the observed
maximum and the subtropical/extratropical totals are low. The single-level atmosphere has no
such equatorial peak.}
\label{fig:precip_zonal}
\end{figure}

\begin{figure}[h]
\centering
\includegraphics[width=\linewidth]{biasmap_precip}
\caption{Precipitation maps: model (left), observed climatology (centre), and their difference
(right). The model captures the broad equatorial rain band and the Pacific/Atlantic structure,
but is systematically too dry along the deep-tropical maximum (brown band near the equator in the
difference map) and lacks the sharp South Pacific and South Atlantic convergence zones.}
\label{fig:precip_bias}
\end{figure}

\section{Assessment}

\paragraph{What works.} The qualitative win is unambiguous and is documented across the model's
development record. Moving from the single-level to the multi-level moist atmosphere lifts the
precipitation \emph{pattern correlation} markedly: in the SST-forced benchmark it rises from
$0.16$ for the single level to $0.34$ for the multi-level atmosphere, and the precipitation
standard-deviation ratio recovers toward observed (from $0.41$ to $\sim0.8$). Crucially, the
correlation gain comes with a \emph{physically correct} ITCZ --- the model now rains on the
equator instead of in the subtropics. In the fully coupled, flux-corrected control run (where the
tropical SST is restored toward WOA so it is warm enough to convect), the same revival is even
clearer: the precipitation pattern correlation reaches $0.37$, and the documented narrative of
that run names the ``revived ITCZ'' as one of its two headline gains. The chain is exactly the
one the physics predicts in Section~\ref{chap:ch19_tropics_monsoon}: warm tropical SST
$\rightarrow$ deep convection $\rightarrow$ Hadley ascent $\rightarrow$ equatorial rain band.

\paragraph{What does not (yet).} Three honest limitations stand out.

\begin{enumerate}[leftmargin=*]
  \item \textbf{The ITCZ is still too dry and too weak.} Figure~\ref{fig:precip_zonal} shows the
    equatorial peak undershooting the observed maximum, and the bias map
    (Figure~\ref{fig:precip_bias}) shows a coherent dry band straddling the equator. A pattern
    correlation of $\sim0.34$--$0.37$ is a real improvement but is still modest. The large-scale
    condensation closure has no convective organization or explicit deep-convection
    parameterization, and at T31 ($\sim3.75^\circ$) the rain band is only a few grid cells wide,
    so it is unavoidably blurred and under-peaked.
  \item \textbf{Tropical SST limits the convection.} The free-running coupled model carries a
    tropical cold-tongue bias (tropical SST several degrees below WOA), which suppresses the deep
    convection that should drive the ITCZ. The control run quoted above relies on a Haney SST
    restoring \citep{haney1971} toward WOA to remove that bias; the ITCZ is therefore partly \emph{earned} by the
    flux correction rather than freely predicted. A free-running warm tropical SST remains a
    standing target (Chapter~\ref{chap:ch13_coupler_fluxes}).
  \item \textbf{No strong interannual (ENSO-like) variability.} At this configuration \model{}
    has no documented El~Ni\~no--Southern Oscillation. The coupled control is run as a stable,
    flux-corrected, drift-free mean state; the ocean's diagnostic-velocity dynamics and the
    coarse T31 atmosphere do not sustain the coupled Bjerknes feedback that produces ENSO. The
    tropical climate documented here is a \emph{mean-state} evaluation, not a variability one ---
    and the spurious noise that the AMOC closure exhibits (Chapter~\ref{chap:ch17_amoc_tipping})
    should not be mistaken for organized tropical variability.
\end{enumerate}

\paragraph{Monsoons.} Monsoon circulations --- the seasonal land--sea thermal contrast that
reverses the low-level winds and drives summer rainfall over South Asia, West Africa, and the
Americas --- are only weakly represented. Two factors limit them. First, resolution: at
$\sim3.75^\circ$ the land--sea coastlines, orographic barriers (the Tibetan Plateau, the
Ethiopian Highlands), and narrow monsoon rain bands are smoothed to a few grid cells, and the
sub-grid Central-American isthmus even has to be closed by hand in the control run
(Chapter~\ref{chap:ch13_coupler_fluxes}). Second, the land surface is a simplified slab plus
bucket-hydrology model (Chapter~\ref{chap:ch12_land}): it captures the first-order land--ocean
temperature contrast but not the soil-moisture--precipitation feedbacks and surface heterogeneity
that sharpen real monsoon onset and intensity. The model produces a continental-scale seasonal
land warming, but the resulting monsoon overturning and rainfall are muted. Sharper monsoons,
like a sharper ITCZ, need higher resolution and a richer land surface.

\paragraph{Why the surface wind pattern lags the rain.} It is worth contrasting the tropical
\emph{thermodynamic} skill (precipitation, which the multi-level atmosphere clearly improves)
with the tropical \emph{dynamic} skill (surface wind and pressure patterns, which it does not yet
improve). The eddy momentum flux that builds the observed surface-wind belts scales with the
equator-to-pole SST gradient; at T31 the regridded WOA SST has only a $\sim22.6$~K gradient,
because the sharp Gulf~Stream, Kuroshio, and ACC fronts are smoothed away. That weak gradient is
enough to organize moisture into an ITCZ but too weak to reproduce the observed surface-wind
pattern --- the same decoupling discussed for the mid-latitudes in
Chapter~\ref{chap:ch18_jets_stormtracks}. Tropical precipitation is the field where the
multi-level atmosphere most clearly earns its added cost.

\section{Summary}

\model{}'s multi-level moist atmosphere produces a genuine ITCZ --- an equatorial rain band
organized by the resolved Hadley circulation --- where the single-level atmosphere structurally
cannot, lifting the precipitation pattern correlation from $\sim0.16$ to $\sim0.34$--$0.37$ and
moving the rain onto the equator. The honest counterweights are that the ITCZ remains too dry and
too weak at T31, that it currently relies on an SST flux correction to keep the tropics warm
enough to convect, that strong interannual (ENSO-like) variability is absent at this
configuration, and that monsoons are limited by resolution and the simplified land surface. The
tropical mean state is a documented success of the multi-level upgrade; tropical variability and
a free-running warm tropics are open work.

% =====================================================================
\section*{Model-as-laboratory exercises}
\addcontentsline{toc}{section}{Model-as-laboratory exercises}

\begin{exercise}[The vertical-structure argument]
Run the single-level atmosphere and the multi-level \dino{} atmosphere on the \emph{same} WOA
SST and compare their zonal-mean precipitation profiles (\code{make\_readme\_figures.py}). Explain,
in terms of the vertical motion field each model can or cannot represent, why only the multi-level
profile has an equatorial peak. Quantify the difference in pattern correlation.
\end{exercise}

\begin{exercise}[SST control on the ITCZ]
The ITCZ should follow warm SST. Force the multi-level atmosphere with a WOA SST field to which
you have added a $+2$~K anomaly in the Northern-Hemisphere tropics and $-2$~K in the
Southern-Hemisphere tropics. Predict, then measure, how the latitude of the zonal-mean
precipitation peak shifts. Relate the shift to the cross-equatorial energy-transport argument of
\citet{vallis2017}.
\end{exercise}

\begin{exercise}[Quantifying the dry bias]
Using the precipitation bias map (Figure~\ref{fig:precip_bias}) from a control checkpoint,
compute the area-weighted mean bias \emph{within} $\pm10^\circ$ of the equator and compare it to
the global-mean bias. Is the deep-tropical dry bias larger or smaller than the global value? What
does this tell you about whether the model's water budget error is concentrated in the ITCZ or
spread evenly? Discuss whether raising horizontal resolution or strengthening the condensation
closure is the more promising fix.
\end{exercise}
